\section{Future Work}

The current study establishes a robust foundational methodology for urban tree detection utilizing Sentinel-2 multi-spectral satellite imagery and a U-Net deep learning architecture. While the developed pipeline successfully addresses several challenges from previous iterations and provides scalable mapping capabilities without relying on costly high-resolution imagery, the exploratory data analysis and implementation constraints reveal several highly promising avenues for future research. This section outlines key areas where the current methodology can be expanded, refined, and integrated into broader environmental modeling ecosystems.

\subsection{Integration with Surface Temperature Modeling and Urban Microclimate Analysis}
One of the most impactful directions for future work lies in the direct integration of our tree detection model with parallel research initiatives, specifically partner projects focused on predicting urban land surface temperatures (LST) and heat islands using Sentinel-2 imagery. Trees play a critical role in cooling urban environments through canopy shading and evapotranspiration processes. Currently, models mapping surface temperatures and models mapping tree coverage operate independently, yet their underlying physical phenomena are inextricably linked. By feeding the high-resolution tree density predictions generated by our U-Net model as a spatial input feature into LST prediction models, researchers could significantly enhance the accuracy of microclimate simulations. Conversely, understanding localized heat anomalies could help identify areas suffering from vegetation stress. Future work should prioritize building a unified, multi-modal pipeline where tree detection and temperature prediction mutually inform each other, ultimately creating a more holistic tool for platforms like \textit{infrared.city} to simulate urban heat mitigation strategies.

\subsection{Exploitation of Extended Spectral Information}
In the current implementation, the model utilizes only four color channels from the Sentinel-2 data stack: red, green, blue, and near-infrared (B02, B03, B04, and B08). This decision was primarily driven by computational limitations and resource constraints during the data preprocessing and model training phases. However, as demonstrated by Freudenberg et al. \cite{freudenberg2025sentinel2}, who developed a Sentinel-2 machine learning dataset specifically for tree species classification, utilizing a broader spectrum of available satellite bands can vastly improve classification accuracy and granularity. Sentinel-2 offers additional channels, such as the Red-Edge bands (B05, B06, B07) and Short-Wave Infrared (SWIR) bands (B11, B12), which are highly sensitive to variations in leaf chlorophyll content, vegetation health, and canopy water stress. Future iterations of this project should expand the input tensor to include these additional bands. Providing the U-Net architecture with a richer spectral signature will likely reduce false positives in dense urban environments and could pave the way for more complex tasks, such as distinguishing broadleaf from coniferous trees or identifying specific dominant species clusters.

\subsection{Resolution of Data Scaling Artifacts and Index Expansion}
During the exploratory data analysis, a systematic data quality issue was identified regarding the calculation of the Enhanced Vegetation Index (EVI) and Soil Adjusted Vegetation Index (SAVI). Because the index formulas were applied directly to the raw pixel values stored in the satellite files---which are large integers---rather than the true light reflectance values they represent, the resulting EVI values falsely exceeded the theoretical upper bound of 1.0 during the summer months. Future work must address this technical debt by strictly dividing the raw band values by 10,000 to yield proper reflectance bounds before index computation. Resolving this artifact will allow the model to safely incorporate both EVI and SAVI alongside the currently used NDVI. Since EVI is less sensitive to atmospheric noise and SAVI corrects for soil background influences, their inclusion is expected to make the model's temporal attention mechanisms much more robust, especially when analyzing autumn or winter imagery where canopy coverage is sparse and background urban surfaces are highly visible.

\subsection{Multi-Sensor Data Fusion to Mitigate Cloud Cover Constraints}
The reliance on optical satellite imagery introduces inherent vulnerabilities to weather conditions. This limitation was distinctly encountered during the data retrieval phase for Paris, where persistent cloud cover in November necessitated substituting the autumn image with data from October. To build a truly resilient and globally scalable system, future work should explore the fusion of optical Sentinel-2 data with Synthetic Aperture Radar (SAR) data from Sentinel-1 \cite{torres2012sentinel}. Radar signals penetrate cloud cover and are highly sensitive to the volumetric and structural characteristics of vegetation (such as trunk and branch geometry). By stacking multi-seasonal SAR features with the existing optical bands, the model could maintain high predictive accuracy even in cloud-prone geographic regions, ensuring uninterrupted temporal monitoring of urban canopies across all seasons.

\subsection{Advanced Spatial Sampling to Counter Imbalanced Label Distributions}
The spatial analysis of the \textit{Baumkataster} datasets revealed significant disparities in urban tree placement patterns across different cities. For instance, while 52.3\% of the sampled trees in Vienna and 42.2\% in Paris are located within designated green OpenStreetMap land use zones, only 2.6\% of Hamburg's sampled trees fall into these areas, indicating a heavy bias toward street-planted trees in the latter. Furthermore, tree density heatmaps demonstrate strong spatial clustering. Training a Convolutional Neural Network on uniformly sampled patches from such spatially imbalanced distributions can lead to models that perform exceptionally well on dense parks but fail to detect isolated street trees. Future methodological improvements should introduce stratified spatial sampling or spatially weighted loss functions (such as Focal Loss \cite{lin2017focal}). This approach would force the network to pay equal attention to sparse urban streetscapes and dense parklands, improving the model's overall generalization and its ability to distinguish trees from surrounding infrastructure without relying solely on general green spaces as a proxy.

\subsection{Scaling Geographic and Climatic Generalization}
Currently, the model's training and evaluation are focused exclusively on cities within the temperate climate zone (Vienna, Paris, Hamburg). While tree cadastre data for Barcelona was collected, it was largely excluded from standard geometry training due to missing size attributes and differing coordinate storage formats. Future work should invest the necessary data-engineering effort to fully integrate Barcelona---and potentially other cities from arid, subtropical, or tropical climates---into the training pipeline. Urban tree phenology in Mediterranean climates, which are dominated by species like \textit{Tipuana tipu} and \textit{Pinus pinea}, differs drastically from temperate zones, exhibiting much flatter seasonal NDVI curves. Expanding the climatic diversity of the training data will be crucial for scaling the application globally, aligning with the ultimate goal of deploying the tool across diverse international contexts to aid global urban climate adaptation efforts.
